\chapter{Testowanie Aplikacji}
\par
\begin{quote}
\"Testowanie jest naturalnym procesem i czy nazwiemy je uczeniem się, czy też eksperymentowaniem, nie ma to większego znaczenia dla samej czynności poznawania i podejmowania decyzji o testowanym obiekcie.\" ~\cite{smilgin2018tester}
\end{quote}
\par
Proces testowania oprogramowania stanowi istotny elemnt cyklu życia aplikacji i pozwala na ocenę poprawności działania system, wykrycie potencjalnych defektów oraz weryfikację zgodności implementacji z założeniami projektowymi.
Testowanie umożliwia również analizę zachowania aplikacji w różnych scenariuszach użytkowania oraz ocenę stabilności i bezpieczeństwa systemu.
\par
W ramach niniejszego rozdziało przeprowadzono analizę działania systemu \gls{Labbyn}.
Proces testowania obejmował weryfikację kluczowych funkcjonalności systemu, poprawności komunikacji z bazą danych, walidacji danych wejściowych, obsługi wyjątków oraz odporności aplikacji na wybranie nieprawdiłowe scenariusze użycia.
Przedstawiono również sodowisko testowe, przykładowe scenariusze testowe, wykryte defekty oraz sposób ich rozwiązania.
\section{Cel i zakres testów}
\section{Środowisko testowe}
\section{Scenariusze testowe}
\section{Cykl życia defektów i proces ich eliminacji}
\section{Wyniki Testów}